\begin{table}[h!]
\centering
\begin{tabular}{| p{0.2\textwidth} | p{0.35\textwidth} | p{0.35\textwidth} |}
\hline
\textbf{Aspek} 
& \textbf{Kondisi Saat Ini} 
& \textbf{Konsep Solusi} \\
\hline

\textbf{Variasi penelitian} 
& \textit{Dataset}, \textit{preprocessing}, dan konfigurasi berbeda sehingga hasil tidak dapat dibandingkan 
& Pipeline seragam untuk melatih dan menguji NB, SVM, dan RF \\
\hline

\textbf{Jenis fitur} 
& Hanya konten teks (TF-IDF) 
& Kombinasi TF-IDF + metadata sumber \\
\hline

\textbf{Metode deteksi di lapangan} 
& Manual, tidak skalabel 
& Pemodelan ML sebagai \textit{benchmark} awal \\
\hline

\textbf{Evaluasi model} 
& Tidak terstandarisasi 
& \textit{Accuracy, precision, recall, dan F1-score} diterapkan dalam dua skenario \\
\hline

\textbf{Kontribusi ilmiah} 
& Pengetahuan terpisah antar penelitian 
& \textit{Benchmark} terukur + analisis pengaruh fitur tambahan \\
\hline

\textbf{Relevansi solusi} 
& Tidak memberi rekomendasi algoritma terbaik 
& Memberi rekomendasi berbasis data untuk konteks hoaks Indonesia \\
\hline

\end{tabular}
\caption{Perbandingan kondisi saat ini dan konsep solusi yang diusulkan}
\label{tbl:perbandingan-kondisi-solusi}
\end{table}
